% subsection: 階差等比複合型 \, $\displaystyle a_{n+1}=pa_n+f(n)$

\subsection{\texorpdfstring{多項式・指数の付加項 \, $\displaystyle a_{n+1}=pa_n+f(n)$}{多項式・指数の付加項}}\label{節：階差等比複合型}
  \begin{bookpoint}
    \textbf{共通の付加項を消す}\par
    付加項が$n$の多項式なら,$a_n$から多項式を引いて等比型にする.付加項が$q^n$なら,$q$の累乗で割る方法も試せる.
  \end{bookpoint}
  \subsubsection{\texorpdfstring{多項式を引く \, $\displaystyle a_{n+1}=pa_n+f(n)$}{多項式を引く}}
    \sectionref*{節：特性方程式型}では,各項から定数を引いて等比型にした.
    今回は付加項$f(n)$が$n$によって変わるため,引く量も$n$によって変えてみる.
    補助数列$b_n=a_n-g(n)$が等比型になるように,引く式$g(n)$を決めよう.
    $f(n)$を多項式とし,$p\ne1$を考えよう.$p=1$なら\bookterm{difference}{階差}型である.

    \bookstep{引く量を選ぶ}
    $b_n=a_n-g(n)$を等比型にするための条件は
    \[
      g(n+1)=pg(n)+f(n)
    \]
    である.この条件なら,
    \[
      \begin{array}{rrl}
        &a_{n+1}={}&pa_n+f(n)\\
        -\,)&g(n+1)={}&pg(n)+f(n)\\ \hline
        &a_{n+1}-g(n+1)={}&p\{a_n-g(n)\}
      \end{array}
    \]
    となり,付加項が消える.
    $a_n=g(n)+b_n$は,\sectionref*{節：非同次と解の組合せ}の「既知の解と\bookterm{homogeneous}{同次}の解」の分解である.
    $g(n)$を探し,$b_n$を\bookterm{initial}{初期条件}から決める.

    まず$a_{n+1}=2a_n+n-1$で試してみよう.
    付加項が一次式なので,引く量も$g(n)=un+v$と置く.
    一次式は添字を一つずらしても一次式のままである.
    元の付加項と次数をそろえ,未知の係数$u,v$で残る項を消そう.
    条件$g(n+1)=2g(n)+n-1$に代入すると,
    \[
      un+u+v=(2u+1)n+2v-1
    \]
    となる.すべての$n\ge1$で成り立つよう係数を合わせると,
    \[
      u=2u+1,\qquad u+v=2v-1
    \]
    より$u=-1,\ v=0$,すなわち$g(n)=-n$を得る.
    従って
    \[
      a_{n+1}+(n+1)=2(a_n+n)
    \]
    となる.初項$a_1+1$から,
    \[
      a_n=(a_1+1)2^{n-1}-n
    \]
    と元へ戻せる.この例は\bookproblemref{practice-basic}{4}{標準問4}でも扱う.

    \bookstep{多項式へ広げる}
    同じ手順を,$f(n)$が$m$次式の場合に広げよう.
    $p\ne1$のときは$m$次の多項式
    \[
      g(n)=x_0+x_1n+\cdots+x_mn^m
    \]
    を候補にする.$g(n+1)-pg(n)$の最高次の係数は$(1-p)x_m$である.
    これを$f(n)$の最高次の係数に合わせ,次に一段低い係数を合わせる.
    $1-p\ne0$なので,この順に$x_m,x_{m-1},\ldots,x_0$を決められる.\footnote{未知の係数を比較して求めるこの方法を\bookterm{undetermined}{未定係数法}という.}

    \bookstep{一般項を求める}
    $b_{n+1}=pb_n$を解き,\,$a_n=b_n+g(n)$へ戻す.
    $p\ne0,1$なら
    \[
      a_n=\{a_1-g(1)\}p^{n-1}+g(n)
    \]
    となる.\footnote{$p=0$なら初項を別に扱い,$n\ge2$で$a_n=f(n-1)$と直接求める.}
    一方,\,$p=1$で$f(n)$が$m$次なら,$g(n+1)-g(n)=f(n)$を満たす多項式$g(n)$は$m+1$次になる.
    例えば$a_{n+1}=a_n+n$では,\,$g(n)=\frac{n(n-1)}2$が同じ更新を満たす.

  \subsubsection{\texorpdfstring{指数で割る  \, $\displaystyle a_{n+1}=pa_n+q^n$}{指数で割る}}
    $a_{n+1}=pa_n+q^n$では,$q$の累乗で割ると付加項が定数になる.$q\ne0$とする.\footnote{$q=0$なら$n\ge1$で$q^n=0$なので等比型になる.}
    付加項$q^n$を$1$にするために,両辺を$q^n$で割る.
    左辺の添字は$n+1$なので,それに合わせて第$n$項を割る量は$q^{n-1}$とする.
    両辺を$q^n$で割り,$b_n=\dfrac{a_n}{q^{n-1}}$と置くと,
    \[
      \begin{aligned}
        \frac{a_{n+1}}{q^n}&=\frac pq\frac{a_n}{q^{n-1}}+1,\\
        b_{n+1}&=\frac pq b_n+1.
      \end{aligned}
    \]
    \booknavref{節：特性方程式型}{定数項を消す方法}へ戻った.
    ただし$p=q$のときは$b_{n+1}=b_n+1$(\bookterm{arithmetic}{等差数列})になり,\bookterm{fixed}{不動点}$\alpha$が存在しない.\footnote{\bookproblemref{example-characteristic-1}{3}{先ほどの例題(3)}と同じ状況である.}
    この場合は$b_n=b_1+(n-1)$として直接求まる.
    いずれの場合も,求めた$b_n$に$q^{n-1}$を掛ければ元の$a_n$に戻る.
    例えば$a_1=1,\,a_{n+1}=2a_n+2^n$(\,$p=q=2$)ならば$b_n=\dfrac{a_n}{2^{n-1}}$が\bookterm{arithmetic}{等差数列}となり,\,$a_n=n\cdot2^{n-1}$が得られる.
